\section{Results}
\label{sec:results}

\subsection{Baseline Digital Performance}
\label{subsec:baseline-digital}

We first establish the digital FP32 baselines for all architectures across the three target datasets. Table 4 summarizes these results. ResNet-18 and Tiny-ViT-5M both perform strongly on CIFAR-10, while Tiny-ViT benefits from transfer learning and reaches the highest accuracy. On CIFAR-100, the gap between the pre-trained Tiny-ViT deployment setting and the from-scratch ConvNeXt training setting widens further, highlighting that the two testbeds serve different but complementary roles: ConvNeXt probes analog adaptation from random initialization, whereas Tiny-ViT probes analog deployment of a compact foundation model. The low-data Flowers-102 regime makes that distinction even sharper: ConvNeXt-Tiny reaches only 33.22\% from scratch, while Tiny-ViT retains 97.97\% after fine-tuning.

| Architecture | Dataset | FP32 Accuracy (%) |
|:--|:--|--:|
| ResNet-18 | CIFAR-10 | 94.98 |
| ConvNeXt-Tiny | CIFAR-10 | 90.74 |
| Tiny-ViT-5M | CIFAR-10 | 97.48 |
| ConvNeXt-Tiny | CIFAR-100 | 64.12 |
| Tiny-ViT-5M | CIFAR-100 | 86.94 |
| ConvNeXt-Tiny | Flowers-102 | 33.22 |
| Tiny-ViT-5M | Flowers-102 | 97.97 |

\subsection{Quantization and the ``Scale Masking'' Effect}
\label{subsec:scale-masking}

Evaluation of the 4-bit hybrid models without explicit noise exposure (Experiment V2/C2) shows that quantization alone is a minor penalty under the canonical mapping. For Tiny-ViT, the more surprising observation is that V2 still attains 97.39\% when evaluated under the standard uniform-noise setting. Our diagnostic interpretation is the \textbf{scale-masking effect}: digital scale recovery compresses the effective analog perturbation below the 4-bit LSB, so the deployment noise is largely absorbed by the quantized weight bins. However, this protection is model-specific rather than unconditional. As Section 5.9 will show, it does not survive under proportional state-dependent noise.

\subsection{Task Complexity Scaling}
\label{subsec:complexity-scaling}

\begin{figure}[t]
    \centering
    \includegraphics[width=0.85\textwidth]{fig4_accuracy_comparison}
    \caption{Cross-dataset accuracy comparison for ConvNeXt-Tiny and Tiny-ViT-5M across three complexity tiers (CIFAR-10, CIFAR-100, and Flowers-102). Tiny-ViT benefits from pre-training but exhibits higher noise sensitivity on complex tasks.}
    \label{fig:accuracy-comparison}
\end{figure}

\begin{figure}[t]
    \centering
    \includegraphics[width=0.85\textwidth]{fig5_hat_recovery}
    \caption{Degradation and recovery summary. Left: Accuracy drop from FP32 to standard noisy evaluation. Right: Recovery gained via Hardware-Aware Training (HAT). Absolute FP32 baselines are indicated to ensure transparency in low-performing regimes like Flowers-102.}
    \label{fig:hat-recovery}
\end{figure}

Figures~\ref{fig:accuracy-comparison} and~\ref{fig:hat-recovery} show the central system-level conclusion of this study: \textbf{noise fragility scales with task complexity, and the value of HAT scales with it.}

On CIFAR-10, the canonical Tiny-ViT degradation from V1 to V3 is modest and HAT provides only a small additional gain. On CIFAR-100, the situation changes sharply: Tiny-ViT drops from 86.94\% to 44.06\% under standard noisy deployment, and HAT recovers the model to 65.48\% (+21.42 pp over V3). ConvNeXt exhibits the same qualitative trend, dropping from 64.12\% to 23.86\% and recovering to 60.54\% (+36.68 pp over C3). This cross-architecture agreement is important because it shows that the complexity-dependent role of HAT is not a transformer-only artifact.

Flowers-102 reveals the opposite boundary condition. Tiny-ViT still benefits from HAT in relative terms (4.81\% to 22.48\%), but remains far below its FP32 baseline of 97.97\%. ConvNeXt does not recover at all (dropping slightly from 3.79\% to 3.35\%). We therefore hypothesize that stochastic hardware-aware regularization has a \textbf{data-volume floor}. Under a shared training recipe and extremely limited supervision, HAT can become too strong for stable convergence, causing the CNN baseline to marginally underperform even the unregularized noisy evaluation. However, this remains a hypothesis that requires further ablation to fully decouple from potential domain-shift factors.

\subsection{ADC and Quantization Thresholds}
\label{subsec:adc-thresholds}

\begin{figure}[t]
    \centering
    \includegraphics[width=0.85\textwidth]{fig9_noise_sensitivity}
    \caption{Continuous noise sensitivity and ADC threshold analysis. The \(6\)-bit ADC knee is a critical practical threshold for Transformer-based CIM architectures.}
    \label{fig:noise-sensitivity}
\end{figure}

The ADC bit-width sweep identifies \textbf{6-bit as the critical practical threshold} for Transformer-based CIM under the present simulator assumptions. At 4-bit ADC resolution, accuracy drops to ~27\%, while a jump to 6-bit restores it to over 80\%. This indicates that the attention mechanism is more constrained by the precision of the summation results than by the static weight programming.

\subsection{Retention and Temporal Drift}
\label{subsec:retention-drift}

\begin{figure}[t]
    \centering
    \includegraphics[width=0.8\textwidth]{fig7_retention_curve}
    \caption{Retention decay for the corrected Tiny-ViT V4 model. The model exhibits a two-phase behavior: a rapid early drop followed by a stable plateau near 79\% accuracy extending out to \(10{,}000\) s.}
    \label{fig:retention-curve}
\end{figure}

Fig.~\ref{fig:retention-curve} shows the corrected V4 retention decay with dynamic scale recalibration and co-decay of the retained D2D buffers. The curve exhibits a two-phase behavior: a rapid early drop from 91.63\% to 82.66\% at 1 s, followed by a broad plateau near \textbf{79\% accuracy} from 10 s to $10{,}000$ s (79.13\%, 79.05\%, 79.35\%, and 79.51\% at 10, 100, 1000, and $10{,}000$ s, respectively). Under the present first-order retention model, this indicates that long-horizon analog inference remains partially viable once gain recalibration is included in the inference stack. The retention results reported here were obtained under the uniform double-exponential decay model (§3). The codebase now additionally supports true state-dependent retention where high-conductance states decay faster, but the canonical V4 results in this paper use the uniform model.

\subsection{Hardware-Instance Transferability}
\label{subsec:hardware-transferability}

\begin{figure}[t]
    \centering
    \includegraphics[width=0.85\textwidth]{fig10_zero_shot_transferability}
    \caption{Zero-shot hardware transferability across different device profiles. Tiny-ViT exhibits severe instance overfitting, while ConvNeXt-Tiny maintains structural robustness across fresh D2D realizations.}
    \label{fig:zero-shot-transfer}
\end{figure}

Experiments on fresh hardware instances (Task 12) reveal a severe weakness in current transformer deployment. Tiny-ViT V4 collapses to near-random performance when evaluated on arrays with a different fixed D2D realization, whereas ConvNeXt retains partial functionality under some transferred device profiles. The key lesson is not that CNNs are universally robust, but that the present Tiny-ViT training procedure remains strongly susceptible to \textbf{hardware-instance overfitting}. This makes same-instance robustness and fresh-instance robustness fundamentally different evaluation targets and motivates multi-instance HAT as a next-step training strategy.

\subsection{Physical Front-end Compensation}
\label{subsec:frontend-compensation}

\begin{figure}[t]
    \centering
    \includegraphics[width=0.8\textwidth]{fig6_physical_compensation}
    \caption{Impact of organic phototransistor non-linearity and dark current. Inverse-gamma compensation provides signal recovery in dark regimes but amplifies shot noise elsewhere.}
    \label{fig:frontend-compensation}
\end{figure}

We evaluate the impact of the organic phototransistor's sub-linear response ($\gamma_{\text{phys}}$) and dark current ($I_{\text{dark}}$) on system-level accuracy. Figure 6 summarizes the ResNet-18 (R4) and Tiny-ViT (V6) results under various frontend configurations.

The inverse-gamma compensation mechanism successfully restores accuracy for high $\gamma_{\text{phys}}$ (e.g., $\gamma=2.0$), providing a gain of up to +5.8 pp compared to the raw physical response. However, we observe that for $\gamma < 1$, the compensation is less effective due to shot-noise amplification. Our SNR analysis (Group 2) confirms that while inverse-gamma preprocessing linearizes the signal, it amplifies noise variance across the entire intensity range when $1/\gamma > 1$.

\textbf{Core Value Conclusion}: Inverse-gamma compensation provides a physical-level regularization effect in low-light/dark scenarios; however, in high-intensity regions, noise variance increases with the compensation exponent, constituting a fundamental \textbf{design trade-off}. Consequently, the frontend benefit is most pronounced in high-dark-current regimes where signal recovery outweighs noise amplification.

\subsection{Physical Non-Idealities}
\label{subsec:non-idealities}

Stress testing the "Scale Masking" and "HAT Recovery" conclusions under more realistic device models reveals substantial risks when models are not explicitly trained for these conditions:

\begin{enumerate}
    \item \textbf{State-Dependent Noise Evaluation}: When a model trained under uniform noise (V4) is evaluated under state-dependent proportional noise, accuracy collapses to \textbf{10.0\%}. This indicates that the uniform window-referenced noise model likely overestimates the intrinsic robustness of cross-device transfer.
    \item \textbf{Nonlinear Writing Evaluation}: Similarly, evaluating V4 under a non-linear writing constraint ($NL=2.0$) during inference-only tests degrades accuracy to \textbf{27.9\%}. This confirms that the global attention mechanism is fundamentally sensitive to weight distortions that deviate from the training-time noise distribution.
\end{enumerate}

\subsection{HAT Recovery under Physical Stress}
\label{subsec:hat-recovery-stress}

To test whether hardware-aware training can overcome these hard physical constraints, we performed dedicated retraining under proportional noise and non-linear writing dynamics for both architectures.

\begin{enumerate}
    \item \textbf{Proportional Noise HAT (Task 34 \& 36)}: Hardware-aware training can recover performance when the training and evaluation physics are aligned. For \textbf{Tiny-ViT-5M}, proportional-noise HAT reaches \textbf{97.48\%} in training and \textbf{97.37 ± 0.05\%} under proportional-noise Monte Carlo evaluation, recovering performance within the matched proportional-noise regime. However, this robustness is highly regime-specific: the same checkpoint falls to \textbf{10.38 ± 0.44\%} when evaluated under the original uniform-noise semantics. For \textbf{ConvNeXt-Tiny}, proportional-noise HAT reaches \textbf{91.98\%} with \textbf{91.91 ± 0.08\%} Monte Carlo performance. Because this slightly exceeds the nominal C1 baseline under the current training recipe, we interpret it as a stochastic-regularization effect under this specific matched setting, not as evidence that analog proportional noise is broadly accuracy-improving.

    \item \textbf{Nonlinear Writing HAT (Task 35)}: In contrast, HAT remains unable to recover performance under severe write non-linearity ($NL=2.0$) for the Transformer backbone. Despite 100 epochs of retraining, Tiny-ViT-5M peaks at \textbf{27.37\%}, and the corresponding evaluation remains at \textbf{27.72 ± 0.82\%}. This is only marginally different from the inference-only degradation, indicating that first-order gradient scaling is not sufficient to bridge the non-linearity gap for the current transformer deployment.

    \item \textbf{Ensemble HAT (Task 37)}: To address the hardware-instance overfitting identified in Section 5.6, we evaluate a "top-tier" model trained with per-epoch D2D resampling. While the standard HAT model (V4) overfits to its training instance and collapses to \textbf{10.00\%} when evaluated on fresh device instantiations, the Ensemble HAT model demonstrates remarkable zero-shot transferability. Across 10 fresh hardware instances with different spatial mismatch maps, the Ensemble HAT model maintains an average accuracy of \textbf{86.37 ± 1.54\%}. This confirms that learning weight basins invariant to specific D2D realizations is a viable algorithmic path to cross-device generalizability for Transformer-class models. While this multi-instance strategy successfully mitigates hardware overfitting, it introduces a measurable training-cost overhead. Because the D2D mismatch mask is resampled every epoch, the forward-pass graph cannot be statically optimized for a single noise realization, which marginally increases the effective training time. However, this one-time offline cost is an acceptable trade-off for achieving zero-shot transferability during inference.
\end{enumerate}

Taken together, these findings revise the physical ranking of bottlenecks. Under the canonical uniform-noise model, standard stochastic device noise is not the dominant limit. Under more realistic extensions, however, the picture becomes regime-dependent: proportional noise can be retrained around, but \textbf{nonlinear write dynamics} and \textbf{instance-specific mismatch} remain the primary challenges for deploying Transformer-class models.

\subsection{Energy Efficiency Profile}
\label{subsec:energy-profile}

\begin{figure}[t]
    \centering
    \includegraphics[width=0.8\textwidth]{fig11_energy_breakdown}
    \caption{Inference energy breakdown for the Tiny-ViT hybrid architecture. Despite analog acceleration of linear layers, digital attention operations remain the dominant energy consumer.}
    \label{fig:energy-breakdown}
\end{figure}

\begin{figure}[t]
    \centering
    \includegraphics[width=0.8\textwidth]{fig8_pareto_energy_accuracy}
    \caption{Pareto-style energy/accuracy summary across various hardware and training configurations.}
    \label{fig:pareto}
\end{figure}

\begin{figure}[t]
    \centering
    \includegraphics[width=0.9\textwidth]{fig_attention_maps}
    \caption{Qualitative attention map comparison across V1, V3, V4, and V6 configurations. HAT effectively restores the spatial attention patterns disrupted by hardware non-idealities.}
    \label{fig:attention-maps}
\end{figure}

The analytical energy model (visualized in Fig. 11) estimates a total inference cost of \textbf{$273.94~\mu$J} for Tiny-ViT-5M, corresponding to an \textbf{11.45x reduction} relative to the FP32 digital reference under the present operation-count assumptions. In the current Pareto view (Fig. 8), this places the hybrid Tiny-ViT regime on a favorable energy-accuracy frontier, but not as a universally optimal point independent of future measured-device calibration or alternative peripheral assumptions. Digital attention operations (\(QK^{\top}\) and Softmax) still account for \textbf{57.9\%} of the total cost, confirming the "analog ceiling" discussed in §6.4.


\subsection{Case Study: Zero-Shot Transfer to a Literature-Calibrated Device}
\label{subsec:case-study}

To demonstrate the practical utility of our framework for evaluating realistic hardware, we conducted a case study simulating zero-shot transfer to a recent state-of-the-art organic optoelectronic synapse array, calibrated directly from literature. We defined a literature-derived profile (`literature\_fitted\_profile.json`) using parameters fitted from the physical device characterization in Zhang et al., \textit{Nature Communications} \textbf{17:197} (2026): a conductance range of $G_{\max}/G_{\min} = 47.3$ with 34 effective programming states, and transparent proxy estimates for programming variability ($\sigma_{c2c}=2\%$, $\sigma_{d2d}=3\%$) derived from the reported supplementary cycling and threshold voltage ($V_{th}$) uniformity data.

By loading this literature-derived organic profile into the simulator, we evaluated the pre-trained Tiny-ViT checkpoints without any device-specific fine-tuning. The standard HAT model (V4), having overfitted to its original training instance, collapsed to \textbf{10.00\%} on the new array parameters. In contrast, the Ensemble HAT model (Task 37) successfully maintained an accuracy of \textbf{88.53\%}. This workflow explicitly confirms that the framework can be employed to benchmark the system-level viability of externally characterized CIM materials, effectively bridging the gap between raw device measurements and algorithm deployment.
